%% =====================================================================
%%  B1-T-11-02 -- what B1 contributes to "The Upper Hand"
%%  -------------------------------------------------------------------
%%  LAYER A: APPEND-ONLY. Written once at the entry's write, then left
%%  alone. If the source has more to say later, add a bullet -- do not
%%  rewrite. Material found ONLY here goes in the `unique` slot with
%%  @unique set to yes, which makes it undeletable (rule R10).
%% =====================================================================
%% @kind: source
%% @id: B1-T-11-02
%% @book: B1
%% @topic: T-11-02
%% @locator: Tactic 4 (ch. V), pp. 43--58
%% @status: distilled
%% @unique: no
%% @integrated: yes
%% @updated: 2026-09-19

\dossier{B1}{T-11-02}{\bookshort{B1}}{Tactic 4 (ch. V), pp. 43--58}

\begin{distilled}
  \item The whole craft is summarised as a state of mind: convey that you do not need the counterpart and they need you.
  \item Perceived need triggers a stubborn reflex --- the target's brakes come on before argument can reach them; perceived independence relaxes them toward the proposal.
  \item The upper hand belongs to whoever can afford to walk away, or can make the other side believe it.
  \item The teacher prepaid the motel room before entering the bar: confidence as infrastructure, rehearsed to the point of never considering failure.
  \item Never the slightest hint of desperation; radiate confidence and independence even when not felt.
\end{distilled}
